\section{Relation to existing literature and exact scope}

The theorem-level conflict is with Campana--Demailly--Peternell
\cite[Theorem~2.2 and Corollary~2.3, p.~680]{CDP2020}.  Their hypotheses match
the now workbench-verified identities
$H^1(X;\mathbb Z)=H^2(X;\mathbb Z)=0$, algebraic dimension one, and
holomorphic algebraic reduction.  Their conclusion gives $c_3(X)\leq0$ and,
for a threefold homeomorphic to $S^6$, algebraic dimension zero.  The same
constructed threefold instead has $c_3(X)=2$ and $a(X)=1$.  It is therefore a
counterexample to Theorem~2.2 and Corollary~2.3 as stated; this conclusion uses
the complete candidate reconstruction, not merely the two printed proof-route
defects below, and is not resolved by observing that the candidate proof does
not invoke \cite{CDP2020}.

The first local issue isolated below occurs in the opening reduction of
\cite[\S7, p.~692]{CDP2020}.  The explicit nodal-cubic family in
\cref{prop:nodal-counterexample} refutes that reduction as a formal
sheaf-theoretic inference.  A second issue is the replacement of monodromy
coinvariants by the full fibre homology in
\cite[Lemma~3.2, pp.~84--85]{CDP1998}; CDP20 later imports the resulting
rank argument.  Either defect might in principle admit a different repair
under stronger hypotheses.  The original theorem cannot be repaired under
its printed hypotheses because the constructed threefold is a verified
counterexample.  The strongest intact replacement is now proved in
\cref{thm:corrected-cdp-global}: one global \(H^2\)-vanishing implies the
stated generic vanishing and \(c_3\le0\).  The stronger relative condition in
\cref{cor:corrected-cdp-relative} is sufficient for such a global vanishing
after a sufficiently positive twist from the base.  For reduced full fibres,
\cref{thm:corrected-cdp-fibrewise} replaces the p.~692 shortcut by the exact
connecting-map injectivity condition.

For the singular fibre, analytic duality is checked in two independent ways.
The proof in \cref{thm:cartier-duality} derives the result from smooth Serre
duality \cite[Theorem~9.37]{Lee2024}; Ramis--Ruget's dualizing-complex theorem
provides the singular analytic comparison
\cite[Theorems~1--2, pp.~81--84, and Proposition~2, p.~91]{RamisRuget1970}.
Relative duality for a proper analytic morphism is also cited in its precise
form as Ramis--Ruget--Verdier, Invent. Math. 13 (1971), 261--283
\cite{RamisRugetVerdier1971}.  Analytic cohomology and base change is
applied in the form of B\u{a}nic\u{a}--St\u{a}n\u{a}\c{s}il\u{a},
Chapter~III, \S3, Theorem~3.4 and Corollary~3.10, together with Grauert's
local-freeness/continuity theorem in Chapter~III, \S4, Theorem~4.12
\cite{BanicaStanasila1976}.  These locators cover all degrees used by the
structure-sheaf direct-image proof, not only its top-degree specialization.
The algebraic top-degree analogue in the Stacks Project, Tag~0EX2
\cite{StacksProject}, is a useful cross-check but is not substituted for the
analytic theorem.

The algebraic nearby-cycle formalism in SGA~7~II, Expos\'e XIII
\cite{SGA7II1973}, and the semistable/Raynaud framework in SGA~7~I,
Expos\'e IX \cite{SGA7I1972}, do not by themselves cover every step for the
present proper holomorphic map, which is not assumed algebraic.  The needed
analytic formalism is nevertheless available.  Dimca's Definition~4.2.1
constructs nearby cycles for complex analytic varieties; Proposition~4.2.2
identifies their stalks with Milnor-fibre cohomology; Proposition~4.2.11
commutes nearby cycles with proper analytic direct image; and
Theorem~2.3.26 gives proper base change, all for coefficient rings including
\(\mathbb Z\) \cite{Dimca2004}.  Kashiwara--Schapira
\cite[\S8.6]{KashiwaraSchapira1990} is an independent formal comparison.
Thus the general integral analytic framework is verified.  None of these
references by itself identifies the candidate's specialization image with
the full monodromy-invariant lattice.  That candidate-specific conclusion,
including saturation, is proved in
\cref{thm:p0-integral-specialization} for the reconstructed proper cusp
model, using the explicit local equations and a supported Gysin calculation.
The direct comparison with the cellular collapse is now proved in
\cref{thm:collapse-equals-specialization}: the neighborhood realization of
the nearby-cycle unit is
\[
 \operatorname{sp}^q=i_F^*(i_W^*)^{-1},
\]
and the candidate retraction gives $(i_W^*)^{-1}=r^*$, hence
$\operatorname{sp}^q=c^*$ over $\mathbb Z$ in every degree.

The topological input to the retraction is Brown's collaring theorem for the
boundary of a topological manifold \cite{Brown1962}, followed by Whitehead and
the cofibration argument in Hatcher \cite{Hatcher2002}.  Clemens's classical
collapse for semistable K\"ahler degenerations \cite{Clemens1977} is useful
context, but it is not used as a substitute for the candidate-specific proof:
the central fibre here is irreducible and nonnormal, so
\cref{lem:collapse-straightening,lem:positive-collapse,%
lem:phase-spread-collapse,prop:verified-cellular-collapse} constructs and
checks the map directly.

For compact-support Poincar\'e duality with local coefficients and its
projection formula we use Bredon \cite[Chapter~VI]{Bredon1997}, with
Hatcher \cite[\S3.H]{Hatcher2002} as a concise topological overview.  This
is the precise input behind the multiplicative identification
\(H^2(B,j_*\mathcal L)\simeq M_G\) used in the Leray calculation.

The completed global comparison uses the chain-level Mayer--Vietoris
convention of Hatcher \cite[\S2.2, pp.~149--150]{Hatcher2002}: the inclusion
of the intersection complex is \(z\mapsto(i_{A*}z,-i_{B*}z)\).  This fixes
the only global minus sign.  The candidate-specific content is not supplied
by that citation: the radial markings in the reconstructed finite fillings
prove that the two band inclusions retract to \(\pi_1\) and \(\pi_2\), while
the explicit toric disc and solid tori give the primitive cusp kernels.
Those geometric facts produce the matrices and the commuting Leray squares
in \cref{prop:actual-finite-band-maps,prop:cusp-mv-blocks,%
thm:mv-leray-comparison}.

For the Hodge-number question, Ugarte's general restrictions on a
hypothetical complex structure on $S^6$ \cite{Ugarte2000} and Angella's
later synthesis \cite{Angella2018} provide the correct general comparison
class.  The candidate-specific calculation is sharper: its reconstructed
structure-sheaf direct images and explicit vanishing of holomorphic forms
first determine the Hodge diamond up to $m=h^{1,2}$ and force
$E_2=E_\infty$.  Neither reference determines $m$ for this family.  The
workbench now does: the cyclic invariant lattices, logarithmic
Kodaira--Spencer divisor, normalized residue sequence, and order-four scalar
give $m=1$, hence $h^{1,1}=2$.  This is a candidate-specific calculation,
not a consequence imported from the general literature.

The spectral interpretation has also been audited against the primary
physical framework.  Jaffe--Witten formulate the four-dimensional
Yang--Mills mass gap on \(\mathbb R^4\) only after construction of the quantum
theory: on its physical Hilbert space, the Hamiltonian \(H\geq0\) has the
vacuum at energy zero and a mass gap precisely when there is a \(\Delta>0\)
such that
\[
  \sigma(H)\cap(0,\Delta)=\varnothing
  \qquad\text{(equivalently, }\sigma(H)\subseteq\{0\}\cup[\Delta,\infty)\text{)}.
\]
This is the physical mass-gap condition in \cite{JaffeWitten2000}.
Euclidean correlation data yield such an
\(\mathcal H_{\mathrm{phys}}\), vacuum, and positive time-translation
generator only through Osterwalder--Schrader reconstruction; the required
input includes Euclidean covariance and symmetry, distributional
regularity/growth, reflection positivity, and clustering, with the precise
reconstruction qualifications supplied by the two primary papers
\cite{OsterwalderSchrader1973,OsterwalderSchrader1975}.  Hodge decompositions,
periods, Leray terms, and coherent direct images are not Schwinger functions
and supply neither reflection positivity nor a physical Hamiltonian.  They
also do not determine even a geometric Laplace spectrum: the explicit
fixed-complex-torus calculation in
\cref{cusp-interpretation-fixed-complex-counterexample} keeps the complex
structure, integral cohomology, Hodge decomposition, coherent direct-image
data, and total volume fixed while the K\"ahler metrics
\[
 g_R=R^2(dx_1^2+dy_1^2)+R^{-2}(dx_2^2+dy_2^2)
\]
have \(\lambda_1(\Delta_{g_R})=4\pi^2/R^2\to0\).  Thus spectra depend on the
chosen metric and operator, and no spectral or physical conclusion follows
from the cohomological/direct-image data alone.  The auxiliary spectral
statements in the workbench are direct operator calculations; none is claimed
to be established by these references or to constitute an
Osterwalder--Schrader reconstruction of quantum Yang--Mills theory.

For the compact-link appendix, the historical scope has now been made
explicit.  Wilson supplies the plaquette/character lattice action
\cite{Wilson_1974}; L\"uscher supplies a selfadjoint strictly positive
transfer-matrix construction in the standard Wilson lattice theory
\cite{Luscher_1977}; and Menotti--Pelissetto establish
Osterwalder--Schrader positivity for the Wilson action in the corresponding
reflection setup \cite{MenottiPelissetto_1987}.  These citations identify the
standard constructions and the physical reconstruction boundary; they do
not prove the appendix's metric-dependent sector normalization, exact flux
action, finite Hessian spectrum, or Gaussian gap, all of which remain the
displayed candidate-specific calculations.  Nor do they turn that finite
Gaussian model into an Osterwalder--Schrader reconstruction, whose axiomatic
requirements remain those of
\cite{OsterwalderSchrader1973,OsterwalderSchrader1975}.

The general mechanisms used inside that candidate-specific calculation were
also checked against the acquired primary archival PDFs (including Deligne's
published erratum); no author- or publisher-supplied primary TeX was found for
these historical works.  Cartan's analytic quotient theorem
identifies the local quotient rings with invariants
\cite[\S1, Theorem~1, pp.~1--2]{Cartan1954Quotient}; Reynolds averaging then
gives exactness of finite-group invariants over $\mathbb C$.  Deligne's
logarithmic conormal and normalized-stratum residue formalism covers the
$d\log$ bases, relative sequence, and self-crossing residue target
\cite[Chapter~II, \S3.2--3.7, pp.~72--78]{Deligne1970}.  Katz--Oda supplies
the connecting-map and exterior-product formalism for the Gauss--Manin and
Kodaira--Spencer maps
\cite[\S2, equations~(5)--(10), pp.~202--204]{KatzOda1968}; Grothendieck and
the cited Stacks tags supply the splitting and elementary-modification
background on $\mathbb P^1$ \cite[Theorem~2.1, p.~126]{Grothendieck1957}.
Finally, Fr\"olicher's original lemmas identify the first and limiting pages
with Dolbeault and associated-graded de Rham cohomology
\cite[p.~643]{Frolicher1955}.  None of these sources supplies the printed
period derivatives or cusp branch maps; those remain the explicit
candidate-specific calculations in the workbench.  The subsequent
line-by-line audit therefore kept their normalization data explicit: the
finite determinant formula is defined up to a nowhere-zero frame unit, and
the constant cusp Kodaira--Spencer matrices are special-fibre reductions
modulo $t$.  These distinctions leave the divisor, residue, and rank
conclusions unchanged.

Once those candidate-specific Hodge and Fr\"olicher data are known, the
Bott--Chern and Aeppli dimensions are no longer open.  McHugh's general
compact-threefold formulas
\cite[Theorem~6.3, Proposition~6.4, equation~(7), and \S7.1]{McHugh2019}
express them in terms
of Dolbeault dimensions, $E_2^{3,1}$, the de Rham subgroup represented by
real $(1,1)$-forms, and the auxiliary spaces $K^{p,q}$.  For this candidate
$E_2^{3,1}=0$, $H^2_{\mathrm{dR}}(X)=0$, and
$\dim K^{p,q}\le b_{p+q}=0$ in every middle degree.  The last inequality is
not an injection $K^{p,q}\hookrightarrow H^{p+q}_{\mathrm{dR}}$: the image in
de Rham cohomology of the pure, $\partial$- and $\bar\partial$-closed forms
surjects onto $K^{p,q}$ because every pure $d$-exact form lies in McHugh's
denominator.  Direct substitution therefore yields the Bott--Chern diamond,
and Schweitzer's Hodge-star isomorphism
\cite[\S2.c, pp.~9--10]{Schweitzer2007} yields the Aeppli diamond, in
\cref{thm:exact-bott-chern-aeppli}.  Angella's published synthesis
\cite[\S4]{Angella2018} independently confirms the endpoint dimensions,
$h^{1,1}_{BC}=2h^{0,1}$, the relation
$h^{1,1}_{BC}+h^{2,2}_{BC}=2h^{2,1}_{BC}+2$, and Schweitzer duality.  The
literature supplies the general identities; the numerical specialization
depends on the workbench's candidate-specific proof that $h^{1,2}=1$ and
$E_2=E_\infty$.

The multiplicative refinement is also now complete.  Schweitzer's
definitions give the Bott--Chern algebra, its natural map to Dolbeault
cohomology, and the Chern-class convention
\cite[\S2 and \S8.a]{Schweitzer2007}, but they do not determine this
candidate's generators or products.  The workbench constructs global
$\bar\partial$-closed representatives $G,Q$, proves their exact marked-fibre
period-coordinate classes, and uses the two exact Fr\"olicher rows, the
factor-$12$ exponential--Leray comparison, Dolbeault--Serre duality, and top
integration to obtain all nine positive-degree Bott--Chern generators and
\[
 H^{*,*}_{BC}(X)=\mathbb C\mathbf1\oplus N,
 \qquad N^2=0.
\]
See
\cref{prop:bcpr-generators-are-bases,thm:bcpr-square-zero-algebra}.  This is
a candidate-specific calculation, not a consequence of the general
dimension identities.  It determines the cohomology classes and every
product, while making no claim of a preferred literal global coefficient
atlas for the noncanonical smooth representatives through the singular
fillings.

The algebraic reduction of a compact complex torus is used in the exact form
given by Birkenhake--Lange \cite[Chapter~2, \S6]{BirkenhakeLange1999}, with
the independent statement in Pareschi--Popa--Schnell
\cite[\S8, p.~2437]{PareschiPopaSchnell2017}.  The toric chart facts are
compared with Fulton \cite{Fulton1993}.  The exact classical-topology formula
\(\pi_1(X_\Sigma)\cong N/N_\Sigma\) is cited from Cox--Little--Schenck
\cite[Theorem~12.1.10]{CoxLittleSchenck2011}, with Danilov's original toric
treatment \cite[\S9]{Danilov1978} as an independent comparison.  All
assertions specific to the locally finite infinite fan and the candidate's
proper-discontinuity argument are proved directly.

Finally, the sphere-recognition endpoint uses Hurewicz and homological
Whitehead exactly as stated by Hatcher
\cite[Theorem~4.32 and Corollary~4.33]{Hatcher2002}, Smale's generalized
Poincar\'e theorem \cite[Theorem~A]{Smale1961}, and the computation of the
smooth endpoint by Smale
\cite[Theorem~1.1 and Corollary~1.2]{Smale1962}.  Kervaire--Milnor provide
an independent equivalent check: their opening table gives
$|\Theta_6|=1$, while the following Remark identifies $\Theta_6$ with
diffeomorphism classes in dimension six
\cite[pp.~504--505]{KervaireMilnor1963}.  Those results close the recognition
step after the candidate's fundamental group and integral homology are
established in \cref{thm:h2-vanishing,prop:geometric-pi1}; they are not used
to establish those inputs.

For the normalized cusp-restriction calculation, the public Stacks Project
references were checked at the formal spectral-sequence level rather than
relying on an unstated edge-map convention.  Leray is Tag~01F2; the passage
$E_r\mapsto E_{r+1}$ and the stable term are Tags~011N and~011O;
functoriality, identification with the associated graded, and convergence
are Tags~012O, 012U, and~012W \cite{StacksProject}.  Together with analytic
proper base change \cite[Theorem~2.3.26]{Dimca2004}, these results make the
evaluation square for
$H^2(X;\mathbb Z)\to H^0(B,R^2f_*\mathbb Z)$ commute with restriction to
every fibre.  Naturality of the exponential-sequence boundary then sends
$c_1(L)$ to $c_1(L|_{X_b})$.  The literature supplies this functorial
formalism; the injective differential
$d_2^{0,2}(2q)=\omega[\delta]$ remains the candidate-specific calculation
proved in \cref{lem:higher-transgressions}.

The independent direct boundary-Picard certificate uses Brown's standard
crossed-homomorphism model for cohomology of a group with coefficients in a
module \cite[Chapter~I, \S5]{Brown1982}.  That reference supplies only the
identification of $H^1$ with crossed homomorphisms modulo principal ones.  The
candidate-specific content is proved directly: the finite-cover norm, the
direction of each logarithmic conjugacy, the contraction formula for a
translated invariant two-form, the order-four parity lattice, the explicit
cusp solid torus, and Chern-class naturality under the local nearby-cycle
specialization.  These give
\[
 E=2nq,\qquad
 (\bar d_1,\bar d_2,\bar d_0)=(4n,-3n,0)[\delta],\qquad n=0.
\]
The exact Appell--Humbert source is Birkenhake--Lange, Chapter~1, \S2,
pp.~5--8, especially Theorem~2.1 on p.~7
\cite{BirkenhakeLange1999}; Chapter~2, \S6 is instead the separate source for
algebraic dimension used above.  Schnell's public notes record the precise
factor/coboundary and antisymmetrized connecting-cocycle formulas
\cite[equations~(3.3)--(3.10), pp.~3--7]{SchnellLecture3}, together with the
translation contraction $w\mapsto\iota_wE$
\cite[equation~(6.2), p.~1]{SchnellLecture6}.  These sources certify the
standard fibrewise formalism, not the candidate's chosen lifts, finite norms,
signs, or cusp winding; those are proved directly.

The completed third certificate uses only holomorphic $\mathcal O^*$ factors,
both finite fillings, the cusp argument principle, and the toric support
function.  Its exact boundary equations are
\[
 E=kq,\qquad
 (\bar d_1,\bar d_2,\bar d_0)
   =(2k,-\tfrac32k,0)[\delta],\qquad k=0.
\]
For normalized central degrees $(d_1,d_2,d_3)$, it then proves
\[
 B=\begin{pmatrix}d_2+d_3&d_3\\d_3&d_1+d_3\end{pmatrix},
 \qquad E_L((\lambda,0),(0,m))=-B(B_0\lambda,m).
\]
Fulton's chart, character-divisor, and support-function results
\cite[\S2.1, pp.~28--29; \S3.3, pp.~60--63; \S3.4,
pp.~63--72]{Fulton1993} and the parallel theorem-numbered statements in
Cox--Little--Schenck
\cite[Example~1.2.21, p.~31; Propositions~4.1.1--4.1.2, p.~171;
Theorems~4.2.8, p.~179 and~4.2.12, p.~182]{CoxLittleSchenck2011}
fix the convention: for $D=\sum a_\rho D_\rho$ the standard support function
has value $-a_\rho$, while
$\operatorname{div}(\chi^m)=\sum\langle m,u_\rho\rangle D_\rho$ has the
positive sign used in the proof.  Cartan theorem~B
\cite[\S2, fundamental theorem, p.~81]{Cartan1957} and the exponential
sequence \cite[Proposition~6.28 and equations~(6.35)--(6.36),
pp.~180--182]{Lee2024} trivialize only the residual line bundles whose
$c_1$ has already been proved zero.  All infinite-fan local finiteness,
descent factors, frame comparisons, and the final cross-block identity remain
candidate-specific proofs in \cref{subsec:analytic-all-picard}.

Bingener's analytic representability theorem was checked against the public
primary text, not inferred from the algebraic Picard literature.  For a
proper flat holomorphic map cohomologically flat in dimension zero, the
relative Picard functor is represented by a complex space which need not be
separated \cite[Introduction, p.~318; Satz~6.2, pp.~336--337]{Bingener1980}.
The same introduction states that the direct exponential-sequence quotient
works under topological local triviality and cohomological flatness in
dimension one, and explains why that quotient argument fails in general.
These statements provide representability for the cusp family but do not
provide the component-specialization or separatedness assertion required to
annihilate the normalized central restriction.  The candidate-specific
calculation in \cref{prop:local-cusp-picard-obstruction} makes this limitation
concrete: it constructs a descended local line bundle with generic Chern
class \(-2q\) and nonzero normalized central degree vector
\((-12,2,0,-12,2,0)\).  Thus all three proved global restriction theorems
genuinely use data absent from a cusp-local representability argument; no
general relative-Picard theorem has silently replaced the candidate-specific
finite-filling and common-boundary calculation.

\subsection{Targeted corpus cross-check (29 August 2026)}

Before extending the validate/repair/disprove chain, the relevant primary
texts were reread in the local LaTeX corpus and compared with their public
editions.  This check adds no unproved hypothesis; it records exactly which
formal tools the sources do, and do not, supply.

First, the primary transcription of Campana--Demailly--Peternell (2020),
Theorem~2.2 and Corollary~2.3, confirms the printed implication
\[
 H^1(X,\mathbb Z)=H^2(X,\mathbb Z)=0,\quad a(X)=1
 \Longrightarrow c_3(X)\leq 0,
 \qquad X\simeq S^6\Longrightarrow a(X)=0,
\]
under the stated compact-threefold and algebraic-reduction hypotheses.  Its
Proposition~2.4 identifies the obstruction with the vanishing of
$R^2f_*(T_X\otimes L)$, equivalently the fibrewise dual spaces involving
$H^0(\Omega_X^1\otimes L^*)$.  The workbench's conductor section,
absolute-fibre duality, and analytic base-change propositions prove a
nonzero singular-fibre contribution for the reconstructed candidate; this is
the exact point at which the earlier conditional objection is replaced by a
candidate-specific calculation.

The same primary source confirms that Proposition~2.4 is logically separate
from the defective attempts to establish its premise.  The workbench now
proves that premise-to-conclusion implication without importing either
defect.  If one \(L\) has
\(R^2f_*(T_X\otimes L)=0\), projection formula and Serre vanishing on the
base produce \(L_0\) with \(H^2(X,T_X\otimes L_0)=0\); generic endpoint
vanishing and semicontinuity then give
\[
 H^i(X,T_X\otimes M)=0\quad(i\ne1),\qquad
 \chi(X,T_X\otimes M)=\frac12c_3(X)\le0.
\]
This is a corrected sufficient theorem, not a reinstatement of Theorem~2.2:
the candidate violates this relative sufficient condition for every \(L\).
Separately, \(c_3(X)=2\) and the corrected global theorem imply by
contraposition that \(H^2(X,T_X\otimes L_0)\ne0\) for every \(L_0\).

The reread of CDP (1998), Lemma~3.2, confirms a separate printed rank defect:
for a split extension with fibre module $H_1(F)$ and monodromy group $Q$, the
abelianization contributes the coinvariants $(H_1(F))_Q$, not the whole
$H_1(F)$.  The elementary example
$\mathbb Z\rtimes_{-1}\mathbb Z$ has abelianization
$\mathbb Z\oplus\mathbb Z/2$, which is the exact correction used in
\cref{prop:monodromy-certificate}.

The SGA~7~II nearby-cycle transcription supplies the specialization--variation
sequence and the identity $\operatorname{Var}(\sigma)\alpha=\sigma-1$;
the Stacks/EGA transcriptions supply Leray functoriality and the precise
proper/flat base-change hypotheses.  Both are algebraic formal references,
so neither is silently applied as an analytic existence theorem.  The
candidate's integral specialization map, saturation, and stalk values remain
the explicit local calculations in \cref{thm:p0-integral-specialization}.

Finally, the independent Angella and McHugh transcriptions reproduce the
known six-sphere Hodge inequalities and the Bott--Chern dimension identities.
Substitution of the workbench values
\[
\begin{aligned}
 (h^{0,1},h^{0,2},h^{1,0},h^{2,0},h^{1,1},h^{1,2})&=(1,0,0,0,2,1),\\
 (h^{1,1}_{BC},h^{1,2}_{BC},h^{2,2}_{BC})&=(2,1,2)
\end{aligned}
\]
agrees with the explicit nine-generator square-zero algebra.  This is an external
consistency check, not a substitute for the candidate-specific quotient,
gluing, and period proofs.

\subsection{Stable-stem audit of the sphere-valued carriers
(3 September 2026)}

Let
\[
 g=\widehat q_3:(S^6,*)\longrightarrow(S^3,*),
 \qquad
 f=\widehat q_\gamma:(S^6,*)\longrightarrow(S^2,*)
\]
be the based sphere-valued quotients obtained after the chosen based
identification of \(X\) with \(S^6\).

Hatcher's full low-dimensional homotopy-group table
\cite[p.~339]{Hatcher2002} gives
\[
 \pi_6(S^3)=\mathbb Z/12,
 \qquad
 \pi_{12}(S^6)=\mathbb Z/2.
\]
In particular \(12g=0\).  Toda's Proposition~5.11
\cite[Proposition~5.11]{Toda1962} identifies the latter group as
\[
 \pi_{12}(S^6)=\mathbb Z/2\{\nu_6\circ\nu_9\}.
\]
The same proposition gives
\(\pi_{n+6}(S^n)=\mathbb Z/2\{\nu_n\circ\nu_{n+3}\}\) for \(n\geq5\).
Suspension carries each displayed generator to the next displayed
generator.  Hence the stabilization homomorphism
\[
 \operatorname{st}:\pi_{12}(S^6)\longrightarrow\pi_6^{\mathrm S}
\]
is an isomorphism.  Isaksen--Wang--Xu's stable-stem table
\cite[Table~1, p.~113]{IsaksenWangXu2023} identifies
\[
 \pi_3^{\mathrm S}=\mathbb Z/24\{\nu\},
 \qquad
 \pi_6^{\mathrm S}=\mathbb Z/2\{\nu^2\}.
\]
Since \(12\operatorname{st}(g)=0\) in
\(\mathbb Z/24\{\nu\}\), there is an \(a\in\mathbb Z/12\) such that
\[
 \operatorname{st}(g)=2a\nu.
\]
Stabilization is multiplicative for smash products, so
\[
 \operatorname{st}(g^{\wedge2})
 =\operatorname{st}(g)^2
 =4a^2\nu^2=0\in\pi_6^{\mathrm S}.
\]
Injectivity of the preceding stabilization isomorphism now gives
\([g^{\wedge2}]=0\in\pi_{12}(S^6)\).  Thus every based map
\(S^6\to S^3\), and in particular \(\widehat q_3\), has based
null-homotopic smash square.  Smashing this based null-homotopy with
further copies of \(g\) also gives
\[
 [g^{\wedge r}]=0\in\pi_{6r}(S^{3r})\qquad(r\geq2).
\]

For the second quotient, Hatcher's stable-range corollary
\cite[Corollary~4.24, p.~360]{Hatcher2002} says that
\(\pi_{n+k}(S^n)\) has reached the stable group \(\pi_k^{\mathrm S}\)
when \(n>k+1\).  Applying this with \(n=8\) and \(k=4\), and then using
the empty \(k=4\) row of Isaksen--Wang--Xu's complete stable-stem table
\cite[Table~1, p.~113]{IsaksenWangXu2023}, gives the exact equality
\[
 \pi_{12}(S^8)\cong\pi_4^{\mathrm S}=0.
\]
Consequently the sixfold reduced suspension
\[
 \Sigma^6f=f\wedge\operatorname{id}_{S^6}:
 S^6\wedge S^6\longrightarrow S^2\wedge S^6
\]
is based null-homotopic.  Under the canonical associativity and symmetry
identifications of based sphere smashes, the smash square factors as
\[
 f^{\wedge2}
 =\bigl(\operatorname{id}_{S^2}\wedge f\bigr)
  \circ\bigl(f\wedge\operatorname{id}_{S^6}\bigr):
 S^6\wedge S^6\longrightarrow S^2\wedge S^2.
\]
Its right-hand factor is \(\Sigma^6f\), so this factorization supplies a
based null-homotopy of \(f^{\wedge2}\).  For every \(r>2\), associativity
gives
\[
 f^{\wedge r}=f^{\wedge2}\wedge f^{\wedge(r-2)},
\]
and smashing the preceding based null-homotopy with the based map
\(f^{\wedge(r-2)}\) proves
\[
 [f^{\wedge r}]=0\in\pi_{6r}(S^{2r})
 \qquad(r\geq2).
\]
Hence both already constructed continuous surjective quotients
\[
 \widehat q_3^{\wedge2}:S^{12}\twoheadrightarrow S^6,
 \qquad
 \widehat q_\gamma^{\wedge4}:S^{24}\twoheadrightarrow S^8
\]
are based null-homotopic.  Neither calculation determines the one-copy
class \([g]\in\pi_6(S^3)\) or \([f]\in\pi_6(S^2)\).  Both arguments use
the sphere-valued postcompositions essentially, so they give no
null-homotopy of any map into a conically stratified target; those based
homotopy classes remain open.  Point-set surjectivity, quotient topology,
the embedded moving four-spheres, and their intrinsic metric limits are
unaffected by nullity of the global sphere-valued maps.

\subsection{Effective counting for the exact Koide sector}

The effective assertion used in the Koide chapter is matched to its literal
domain rather than transferred from a smoother or radial problem.  After the
five defining walls and the height \(H(a,q,n)=4a+17n\) are checked directly
against Definition~13.1, Benoist--Oh, Theorem~13.2
\cite{BenoistOh_2012}, applies to each of the finitely many integral orbits
and proves
\[
 R(X)=C_0X+O(X^{1-\eta})
 \qquad(\eta>0\text{ fixed but unspecified}).
\]
Neither the statement nor the proof of that theorem supplies
\(\eta\geq\tfrac12\).

The hyperbolic literature gives two stronger exponents for different
domains.  Liu--Sarnak, Proposition~3.2 \cite{LiuSarnak_2010}, gives the
smooth standard-radial exponent \(1/2+7/64=39/64\).  The exact order group
is a congruence arithmetic Fuchsian group because, for a maximal overorder
\(\mathcal M\),
\[
 \mathcal M^1(18\mathcal M)\subset\mathcal O^1\subset\mathcal M^1;
\]
Jacquet--Langlands and Kim--Sarnak give
\(\lambda_1\geq975/4096\) for that input
\cite{KimSarnak_2003,ChuLi_2016}.  Good's sharp radial theorem and its
Chatzakos--Petridis formulation \cite{Good_1983,ChatzakosPetridis_2016}
instead have a \(2/3\) error after the small-spectrum terms are displayed.
The exact radial Huber variable is \(7a+15\sqrt2\,n\), whereas the Koide
count uses \(4a+17n\) together with five sharp walls.  Consequently neither
radial theorem transfers its exponent to the target count.  A pointwise
square-root remainder for the unchanged domain remains open; the complete
hypothesis-by-hypothesis audit is \cref{koide-lit-determination}.

\subsection{Constrained dynamics and symplectic reduction}

The worldline calculation in
\cref{wpr:sec:worldline-reduction} uses three primary sources for its
general reduction framework.  Dirac's original constrained-Hamiltonian
paper \cite{Dirac1950} supplies the distinction between regular and singular
constraint dynamics and motivates retaining the multiplier and all primary
constraint equations.  Gotay--Nester--Hinds
\cite{GotayNesterHinds1978} supplies the presymplectic constraint algorithm:
one first restricts the closed two-form to the constraint locus, determines
its characteristic distribution, and checks whether successive consistency
conditions or rank jumps occur.  Marsden--Weinstein
\cite{MarsdenWeinstein1974} supplies the regular symplectic quotient by a
Hamiltonian symmetry action.

Their role here is structural, not candidate-specific evidence.  None of
these sources computes the Gram morphism
$P_\alpha(Z)=\alpha Z^*Z$, its $SU(3)$-orbits and stabilizers, the
rank-one zero of $\mathrm d(\det\circ P_\alpha)$, the apex branch, or the
hard-constrained packet's threshold rank jump.  Those data and the exact
cotangent sign are established directly in the cited workbench module.  In
particular, regular Marsden--Weinstein reduction applies to the massive
rank-two stratum and to the intrinsic punctured rank-one stratum after their
characteristic quotients have been proved regular; it does not authorize an
extension of that quotient through the irregular zero-momentum apex.

\subsection*{Quaternionic linear algebra for the 4 September continuation}
The local primary corpus was searched for quaternionic determinant and
Jordan-cone sources before promoting the two new Pro packets. The available
Alesker paper, \emph{Non-commutative linear algebra and plurisubharmonic
functions of quaternionic variables}, arXiv:math/0104209, Theorem 1.1.9,
gives the Moore determinant restriction and congruence law.
Wei Wang, \emph{The linear algebra in the quaternionic pluripotential theory},
arXiv:1805.07759v1, Propositions 2.1--2.3, supplies the corresponding
spectral theorem and complex matrix representation. Its primary LaTeX
source was retrieved and inspected. The sign conversion from its
representation to the present one is conjugation by
$\operatorname{diag}(I,-I)$, stated explicitly in
Section~\ref{sec:q444-reconciliation}.
These classical results are cited as background. The marked thirteen-slice
identity, nonunitary product coefficient, and scalar normal-force witness
are calculated explicitly here rather than attributed to those papers.

For the scalar-to-gauge comparison, David Tong's \emph{Lectures on Gauge
Theory}, Section 2.2.1, equation (2.27), identifies the $A_0$ equation
that must be retained after choosing temporal gauge:
\url{https://www.damtp.cam.ac.uk/user/tong/gaugetheory/gt.pdf}.
The primary PDF was obtained and the surrounding derivation read; the
local corpus was first checked for the corresponding moment-map and
constraint references. This workbench independently varies its own
anti-Hermitian $A_0$ and derives the matter moment
$[\Phi,\dot\Phi]+b\dot b^*-\dot b b^*$, so no coefficient or sign is
transferred from Tong's distinct field convention. The exact moving-frame
transformation retains the resulting nonzero Gauss residual.

For the vacuum-orbit comparison, the primary LaTeX of Andrea Gambioli,
\emph{Eight-dimensional $SU(3)$-manifolds of cohomogeneity one},
arXiv:math/0611796v1 (27 November 2006), was retrieved and its
$\mathbb{HP}^2$ example read:
\url{https://arxiv.org/abs/math/0611796}.
This is classical cohomogeneity-one geometry, not a new classification
claimed here. The displayed family and its invariant $\|b\|^2$ are
verified directly in Proposition~\ref{q444:vacuum-average}; the exact
bridge proved there is the Haar average of the marked Lorentz determinant.
The source archive and extracted TeX SHA-256 values are retained in the
supplement's \path{GAUGE_METRIC_SOURCE_RECEIPT.md}, together with the
separate Tong source used to check retention of Gauss' equation.

\subsection*{4 September: source roles in the integrated kinetic and metric continuation}
The canonical quaternionic proof now also uses the classical
gauge-invariant vector-operator context in Tong, \emph{Lectures on Gauge
Theory}, Section~2.7.3, printed pages~107--108,
\url{https://www.damtp.cam.ac.uk/user/tong/gaugetheory/gt.pdf}.
The model-specific eight-eigenvalue inertia calculation, all transverse
field equations including Gauss, conserved current, coefficient
$-4\|b\|^2$, and relation $\mu_{\mathrm d}^2=(2g^2/3)\|b\|^2$
are proved directly in Theorem~\ref{q444:gauge-spectrum} and
Proposition~\ref{q444:current-mass}; they are not imported from Tong's
different scalar example. Gambioli's Section~4 remains background for the
orbit types only. The original kinetic norm, nonzero normal acceleration,
exact current-free periodic continuation, and null waves are independently
proved in Propositions~\ref{q444:vk-horizontal}, \ref{q444:vk-sector},
Theorem~\ref{q444:vk-periodic}, and Proposition~\ref{q444:vk-wave}.

The flat-torus conventions recorded in the integrated metric section are
A.~Gogolev, \emph{Spectra of flat tori}, Section~2 (27 November 2017),
\url{https://andreygogolev.wordpress.com/2017/11/27/spectra-of-flat-tori/},
and V.~Ivrii, \emph{Partial Differential Equations}, Appendix~4.B,
\url{https://www.math.utoronto.ca/courses/apm346h1/20181/PDE-textbook/Chapter4/S4.B.html}.
The present integer-valued dual is $\mathcal N^{-\mathsf T}\mathbb Z^6$;
Ivrii incorporates $2\pi$ in the dual instead. U.~Lang,
\emph{Lecture Notes on Riemannian Geometry}, 16 June 2020,
Definition~4.1 and Proposition~4.4, printed pages~45 and~47,
\url{https://metaphor.ethz.ch/x/2020/fs/401-3532-08L/sc/DG2_16June2020.pdf},
supplies the standard horizontal-isometry definition of a Riemannian
submersion. The specialized proofs remain in
Theorems~\ref{q444:hermitian-cometric}--\ref{q444:metric-submersion}
and Proposition~\ref{q444:native-period-spectrum}: $W$ is a lattice
cometric, $h=2H^{-1}$, native $h=I$ gives $H=2I$, and the original
source/quotient lattice gaps have leading constants $4\pi^2$ and
$8\pi^2$ over $y^2$. These are source-convention pointers to the
integrated proofs, not claims of a new independent source retrieval in
this metadata pass or of a physical quantum spectral conclusion.

\paragraph{4 September: finite homogeneous quantization.}
The complete proof in Section~\ref{qh:chapter} uses the classical
bosonic-matrix mechanism of B.~Simon,
\emph{Some quantum operators with discrete spectrum but classically
continuous spectrum}, Annals of Physics \textbf{146} (1983), 209--220,
\url{https://doi.org/10.1016/0003-4916(83)90057-X}.
The original author-hosted scan was read, especially Section~2 and
Section~7, equation~(3) and Corollary~4. The present proof gives a direct
slice-oscillator and finite-cell argument in the original
$-\operatorname{tr}$ metric, rather than invoking a matrix-model
discreteness theorem. The root-frequency constant is $2\sqrt6$,
including eigenvalue walls.

The bounded-below Brownian-bridge formula is Theorem~1.1 of
B.~Simon, \emph{A Feynman--Kac Formula for Unbounded Semigroups},
arXiv:math-ph/9907022v1 (1999),
\url{https://arxiv.org/abs/math-ph/9907022}.
Its original LaTeX source was retrieved and read. The proof here
retains the anisotropic diffusion matrix and mass-volume measure.
D.~Tong, \emph{Lectures on Gauge Theory}, Section~2.2.1,
equations~(2.27)--(2.31),
\url{https://www.damtp.cam.ac.uk/user/tong/gaugetheory/gt.pdf},
supplies the Gauss-after-temporal-gauge convention. All triplet,
adjoint and volume factors are derived here, not copied from a
different gauge normalization. The corresponding infinite-volume
and continuum interacting spectral question remains open.

\paragraph{4 September: moving Fourier-state comparison and retained Gauss law.}
The original TeX of G.~Panati, H.~Spohn and S.~Teufel,
\emph{The time-dependent Born--Oppenheimer approximation},
ESAIM: M2AN \textbf{41} (2007), 297--314,
\url{https://arxiv.org/abs/0712.4369},
\url{https://doi.org/10.1051/m2an:2007023}, was retrieved and read
for the projected Hamiltonian, Berry and Born--Huang terms, and the
Duhamel/invariance distinction in Sections 1--2.1.
In particular compression alone does not justify a dynamical
approximation or higher-order expansion. Section~\ref{fq:section}
therefore proves an explicit residual and fixed-physical-time estimate
directly for the retained quartic two-mode operator; no general
Born--Oppenheimer theorem is asserted as new or invoked with unchecked
hypotheses. The positive ground-state and parameter regularity
arguments are included. Tong's cited Section 2.2.1 was also reread:
fixing temporal gauge does not remove Gauss. The complex rotating-pair
proof above retains it with the original anti-Hermitian conventions.

\paragraph{5 September: exact coupling elimination and spatial projection.}
The Sylvester map in
\eqref{q444:generic-sylvester-map} is the classical polynomial map
$(A,B)\mapsto AF+BG$; the primary Stacks Project text,
\url{https://stacks.math.columbia.edu/tag/00UA}, was checked for its
matrix description and coprimality criterion. The proof here includes
the degree-drop argument, the explicit retained coefficient ring,
the nonzero determinant witness and the elementary measure-zero proof.
No resultant theorem is presented as new. The new applications are
the exact degree-twenty-eight residual polynomial, its complete
coupling-fibre inverse, and the nonzero spatial-projection polynomial
for the specified rotating gauge ansatz. All metric, Gauss and scalar
conventions are derived in Section~\ref{sec:q444-reconciliation};
they are not inferred from a different model or an assumed quantum
spectral comparison.

\paragraph{5 September: two-cross residual certificate.}
The local primary corpus was checked, including David Tong's
\emph{Gauge Theory}, printed p.~40, equations (2.27)--(2.28),
\url{https://www.damtp.cam.ac.uk/user/tong/gaugetheory/gt.pdf},
on retaining the temporal Gauss constraint. The full matter-current
equations used here are derived in the retained-system subsection, not
imported from the pure-gauge example. The primary Stacks Project
Example 10.143.12, Tag 00UA, supplies the classical Sylvester-map
context. The degree-drop kernel, block factor, positive residual identity
and Cauchy--Schwarz constants in Section~\ref{q444:twocross:section}
are proved directly for the literal data. Neither standard linear algebra
nor Gauss law is represented as a new discovery or a quantum conclusion.

\paragraph{5 September: the modulo-$840$ lattice and its full Weyl quotient.}
Cheng--Duncan--Harvey \cite[Table~3, arXiv p.~21]{esn:CDH2014}
already record the group $GL_2(3)$ for the $A_5^4D_4$ Niemeier lattice;
their equation (2.30) defines the umbral group as the automorphism group
modulo the Weyl group. This abstract identification is not claimed as new.
Section~\ref{esn:sec:shared} proves the particular residue-indexed
index-$72$ glue, its exact minimum table, the transporter from the earlier
glue, and the signed code isomorphism with $GL_2(\mathbb F_3)$.
Theorems~\ref{esn:fa:thm:full-quotient}
and~\ref{esn:fa:thm:residue-weyl} exhaust all lattice isometries as
$wL_g$ and place the literal residue six-cycle in the Weyl group,
centralizing this chosen complement, with every fibre origin retained.

Mare--Willems \cite[Definition~2.2, Propositions~2.4 and~2.7,
and equations (5.4)--(5.6)]{esn:MareWillems2013} supply the flag
$F_4/\operatorname{Spin}(8)$ and its vector/even-half-spin/odd-half-spin
tangent modules. The marked dual classes $v,s,c$ are matched to their
shortest weight sets by the coordinate proof in
Theorem~\ref{esn:thm:d4-f4-triality}. The subsequent exact isometry
in Theorem~\ref{n24:thm:full-isometry} maps the entire rank-$24$
real span into the original three octonion sectors, with inverse,
glue test and the specified cyclic intertwining equation.
The retained residual and Foundation sources
\cite{esn:Residual2026,esn:ESMain2026,esn:Provenance} provide the earlier
shell and glue data. Their exact excerpts and snapshot digests accompany
the source package. The raw count and complementary-divisor quotient
are separately proved, with $A_3(2521)=0$, $A_3(3361)=6$ despite equal
prime residues modulo $840$. The Fourier map retains the actual finite
prime sample and its counts; neither universal selector positivity
nor a physical model is asserted by these lattice and counting maps.

\paragraph{6 September: the Leech and full triality comparisons.}
Conway--Sloane \cite[pp.~277--280]{esn:ConwaySloane1982} give the
holy coefficient construction. Borcherds \cite[Sections~6--7]{esn:Borcherds1985}
proves uniqueness and the uniform construction theorem.
Theorem~\ref{esnl:thm:leech-neighbor} now writes that construction for
the actual $72$ marked glue words: $\rho^2=84$, the two index-six
inclusions, all shifted-class minima and every norm-two exclusion are
proved in the original Euclidean coordinates. The original
$F_4$ Coxeter matrix also records its full $D_4$ discriminant action,
including $Cv_D=c_D+(1,-1,0,0)$, not equality of the chosen representatives.
Yokota \cite[Sections~1.10, 1.14, 1.16; Theorem~2.7.1]{h24:Yokota2009}
fixes the related-triple and Albert conventions. The workbench proves
the complete Peirce table, the determinant coefficients, exact Clifford
maps, and global associated-bundle descent in those conventions.
Mare--Willems \cite[Appendix~B, p.~519; orientation convention pp.~495--496]{esn:MareWillems2013}
supplies the real-flag cell theorem. The explicit finite exponential,
inverse chart, orientation sign, quotient continuity and homology maps
are all computed here. The Hopf pullback and triality reduction use
the previously stated retained $H$ and based source identification;
they do not silently recertify the global complex-threefold construction.
